\scanpage{356}
% top centre
\mspage{272}
% line 1
foncteur canonique pl. fidèle
% line 2
\[
\operatorname{Ind}_\pi(\underline{\operatorname{Hom}}(J,C))\longrightarrow\underline{\operatorname{Hom}}(J,M)
\]
% lines 3–5
déduit de l'inclusion $C\unsure{\mathrel{\lhook\joinrel\relbar}|\hookrightarrow} M$ \deleted{est} d'où $\underline{\operatorname{Hom}}(J,C)\unsure{\mathrel{\lhook\joinrel\relbar}|\subset}\underline{\operatorname{Hom}}(J,M)$ (cf. 4.9) est une équivalence de catégories.

% lines 6–7
(\uline{NB}. J'ignore s'il en est encore ainsi, sans hyp de rigidité sur $J$.)

% lines 8–11
\uline{Démonstration} Cela \deleted{résulte du} \inserted{se prouve comme le} résultat \deleted{\unsure{coroll|coroll.}} similaire pour $\operatorname{Ind}(C)$ \unsure{ald|a l d} $\operatorname{Ind}_\pi C$, \deleted{\unsure{prouvé dans|prouvé ds}} \inserted{cf} SGA 4 I 8.8.5, \deleted{\unsure{et div|et dans}}

% lines 12–23; diagonally cancelled block
\deleted{\uline{Corollaire} 4.11.7. Soient $C,\pi,M$ comme dans 4.11.6, et soit $J$ une catégorie \deleted{telle que $\unsure{\operatorname{card}\,\mathrm{red}\,J|\operatorname{card}\,\mathrm{réd}\,J|\operatorname{card}\operatorname{Fl} J|\operatorname{card} J}\leq\pi$\unsure{)|}} \inserted{équivalente à une catégorie finie,} enfin soit $\varphi\colon J\longrightarrow M$ un foncteur. Pour que $\varphi$ admette une ``$\pi$-représentation indicielle'', i.e. pour que $\varphi$ soit dans l'image \inserted{essentielle} de $\operatorname{Ind}_\pi\underline{\operatorname{Hom}}(J,C)\longrightarrow\underline{\operatorname{Hom}}(J,M)$, il (\unsure{f. et|f. et.}) suffit que $\varphi$ admette une représentation indicielle sans plus, i.e. que son image dans $\underline{\operatorname{Hom}}(J,\operatorname{Ind}(C))$ soit dans l'image essentielle de $\underline{\operatorname{Ind}}(\underline{\operatorname{Hom}}(J,C))$.}
% archive sheet number: 355
